% Figure 14: The epsilon-neighborhood of M
\begin{tikzpicture}[line cap=round,line join=round,>=Stealth,scale=0.9]
  % tube surface
  \draw[thick] (-2.7,0.6) .. controls (-1.1,1.1) and (1.1,1.1) .. (2.7,0.5);
  \draw[thick] (-2.7,-0.9) .. controls (-1.1,-0.4) and (1.1,-0.4) .. (2.7,-1.0);
  \draw[thick] (-2.7,-0.15) ellipse (0.22 and 0.75);
  \draw[thick] (2.7,-0.25) ellipse (0.22 and 0.75);
  % core curve M, dashed inside the tube
  \draw[thick] (-3.5,-0.05) -- (-2.7,-0.15);
  \draw[thick,dashed] (-2.7,-0.15) .. controls (-1.2,0.3) and (1.2,0.3) .. (2.7,-0.25);
  \draw[thick] (2.7,-0.25) -- (3.5,-0.35);
  \node at (-3.45,0.28) {$M$};
  % r(x), x, and the outward normal g(x)
  \fill (0,0.18) circle (1.5pt);
  \node[below] at (0.12,0.1) {\footnotesize $r(x)$};
  \draw[dashed] (0,0.18) -- (0,0.96);
  \fill (0,0.96) circle (1.5pt);
  \node[left] at (-0.08,1.02) {\footnotesize $x$};
  \draw[->,thick] (0,0.96) -- (0,1.65);
  \node[right] at (0.08,1.45) {$g(x)$};
  \node at (2.45,1.0) {$N_\epsilon$};
\end{tikzpicture}
